\chapter{创新性说明}
\section{把来源用途落实到跨角色和跨会话执行}
作品把来源记录与工具授权相连：父来源和用途约束随摘要、委派和受控记忆传递，执行时再结合任务权限检查。具体实现包括子Grant收窄、记忆版本与CAS、来源恢复，以及结构化返回的完整快照和文本字段定位。其工程价值是让一次动作能追溯到输入资料、任务权限和实际结果；当前证据覆盖指定委派与记忆实例。

\section{分别核验文本风险与动作授权}
规则初筛回答“这段资料是否含有风险信号”，契约与策略回答“这个动作是否被授权”。两类结果独立留痕，敏感写入目标另受不可覆盖名单保护。这样，即使一段文本漏检，也仍可核查动作因何被允许或拒绝。

冻结版\texttt{rules-v0.3}由历史协议固定源文件哈希，扩展版\texttt{rules-v0.6}承载后续规则。版本选择与样例哈希共同记录，便于解释同一资料在不同规则版本下的结果，避免规则升级静默改变历史实验。

\section{以账本事实解释违规链路}
溯源图把来源构件、主体、工具调用、决策和影响目标连接起来，保留重复使用与跨角色传播关系。操作者可从阻断动作回查输入资料，也可核对批准对象与实际收据。图中的上下文关联用于定位调查线索；是否发生危险效果仍由执行与状态记录确认。

\section{把业务保全纳入安全验收}
危险原文在安全研判中也是业务证据。系统允许受限分析和字面引用，将业务完成与危险效果分开登记。历史比较保留两个对照配置的实际危险推送和三个标记写入，未用“运行正常”替代“安全完成”。这种验收方式能检查防护是否在拒绝越权操作的同时保留了任务所需资料。

\section{绑定审批对象与实际提交}
审批关联确切输入、版本、目的地和资源状态，提交时再次复核。IOC与Git流程将批准和执行分为两步，对漂移、拒绝、重放及UNKNOWN分别留痕。它为“批准了什么、实际发生了什么”提供可核对的对象，目前保证范围限于进程内状态。

\section{与已有机制的关系}
表\ref{tab:compare}列出相关工作的主要思路及本作品对应实现。该表用于说明设计来源与工程侧重点，不是防护强度排名，也不声称本作品在所有维度优于已有方法。

\begin{table}[htbp]\centering
\caption{相关工作与本作品的工程侧重点}\label{tab:compare}
\begin{tabularx}{\textwidth}{P{0.22\textwidth}YY}\toprule
相关工作 & 公开描述的主要思路 & 本作品对应实现\\\midrule
CaMeL\cite{camel} & 分离控制与数据流，以能力约束信息流 & 来源用途信封与执行时权限检查\\
Task Shield\cite{taskshield} & 核验指令和工具动作是否服务用户任务 & 任务契约；阻断后的对齐说明\\
AgentSpec\cite{agentspec} & 以DSL描述并执行运行时约束 & 受限DSL、效果适配器与三态决策\\
AuthGraph\cite{authgraph} & 比较授权图与执行来源图 & 事件账本的离线溯源展示；未实现同等双图判定\\
AgentVisor\cite{agentvisor} & 语义权限分离与违规反馈恢复 & 可信宿主门禁、反馈与有界循环\\\bottomrule
\end{tabularx}
\end{table}

ContainmentBench关注暴露后的遏制评估\cite{containment}，AID-Guard关注委派效果从授权到提交、重试与恢复的状态化约束\cite{aidguard}，MCP-SafetyBench与MCPTox提供MCP安全基准与工具投毒研究\cite{mcpsafety,mcptox}。本作品当前贡献集中在上述机制的可运行组合、跨角色与跨会话对象，以及独立业务/副作用核验；尚未通过统一真实模型实验测量相对收益。
